\section{Introduction}
\label{sec:intro}

An autonomous agent with finite onboard resources repeatedly faces a consequential stopping problem: when should it abandon productive work and commit to replenishment? A conservative decision wastes usable capacity and lowers task throughput; a late decision can strand the agent before it reaches a charger. This tension appears in delivery robots, mobile manipulators, field vehicles, and aerial systems. It is not resolved by optimizing nominal task reward alone because the return decision is one-way within a resource cycle and its cost is asymmetric.

\begin{center}
  \centering
  \resizebox{0.98\linewidth}{!}{\begin{tikzpicture}[
  x=1cm,y=1cm,
  font=\scriptsize,
  box/.style={draw=#1!85!black,fill=#1!8,rounded corners=2pt,minimum height=0.56cm,align=center,inner xsep=4pt},
  arr/.style={-{Latex[length=2.0mm]},thick,draw=#1},
  note/.style={align=center,font=\scriptsize},
]
\definecolor{oiBlue}{HTML}{0072B2}
\definecolor{oiOrange}{HTML}{E69F00}
\definecolor{oiGreen}{HTML}{009E73}
\definecolor{oiVerm}{HTML}{D55E00}

\node[box=oiBlue,minimum width=1.25cm] (task) at (0,0) {task stream};
\node[box=oiBlue,minimum width=1.15cm] (pi) at (1.55,0) {policy $\pi$};
\node[box=oiOrange,minimum width=1.40cm] (filter) at (3.25,0) {hard safety\\operator $\Pi$};
\node[box=oiGreen,minimum width=1.55cm] (plant) at (5.20,0) {plant/resource\\primitive $K$};
\node[box=oiGreen,minimum width=1.45cm] (rtg) at (7.25,0) {Resource-to-Go\\$\widehat{\mathcal E}_C$};
\node[box=oiBlue,minimum width=1.35cm] (rel) at (9.15,0) {epistemic\\reliability};
\node[box=black,minimum width=1.85cm] (manager) at (11.35,0) {one-way\\ReturnManager};

\draw[arr=oiBlue] (task) -- (pi);
\draw[arr=oiBlue] (pi) -- node[above,font=\tiny] {$a^{\rm nom}$} (filter);
\draw[arr=oiOrange] (filter) -- node[above,font=\tiny] {$u^{\rm exec}$} (plant);
\draw[arr=oiGreen] (plant) -- (rtg);
\draw[arr=oiBlue] (rtg) -- (rel);
\draw[arr=black] (rel) -- (manager);

\node[box=oiGreen,minimum width=1.20cm] (cont) at (10.25,-0.90) {\textsc{Continue}};
\node[box=oiVerm,minimum width=1.55cm] (ret) at (12.20,-0.90) {\textsc{Commit to charger}};
\draw[arr=oiGreen] (manager) -- (cont);
\draw[arr=oiVerm] (manager) -- (ret);

\node[box=oiOrange,text width=2.45cm,minimum height=0.72cm] (early) at (1.45,-1.02) {\textbf{too early}\\unused resource $\Rightarrow$ lower throughput};
\node[box=oiVerm,text width=2.65cm,minimum height=0.72cm] (late) at (4.60,-1.02) {\textbf{too late}\\insufficient return resource $\Rightarrow$ stranding};
\node[note,anchor=north west,text width=6.45cm] at (-0.05,-1.60) {\textbf{Scientific chain:} pair shift $\to$ executed-interface extrapolation $\to$ long-horizon error $\to$ boundary error};
\node[note,anchor=north west,text width=5.65cm] at (7.15,-1.60) {\textbf{Safety scope:} HOCBF handles collision constraints; the learned resource module is not a certificate.};
\end{tikzpicture}
}
  \captionof{figure}{\textbf{Return decisions depend on the executed closed loop.} The task policy proposes \(a^{\rm nom}\), the hard safety operator produces \(u^{\rm exec}\), and the shared plant/resource primitive determines charger Resource-to-Go. The learned module estimates this requirement and may abstain when epistemic reliability is poor; it does not certify collision safety. Committing too early reduces throughput, while committing after the available-resource boundary risks stranding.}
  \label{fig:overview}
\end{center}

State of charge and charger distance are tempting decision variables, but neither specifies the resource required by the trajectory that will actually be executed. A learned policy proposes a nominal action, a safety operator may replace that action, and the plant then realizes motion and resource use. Thus the relevant chain is
\(
\pi\!\to\! a^{\mathrm{nom}}\!\to\!\Pi\!\to\!u^{\mathrm{exec}}\!\to\!K\!\to\!\RtoGo
\),
where \(\RtoGo\) is charger Resource-to-Go. An obstacle that activates the filter can make two equally distant states require different maneuvers and resources. Figure~\ref{fig:overview} states the resulting decision problem and the scope of each component.

Existing work addresses parts of this chain. Trajectory-energy models estimate consumption and risk \citep{choudhry2021cvar}, and predictive resource-safety networks place a learned override above task policies \citep{guo2020predictive}; constrained and distributional RL optimize policies under resource, expected-cost, or risk-sensitive constraints \citep{qin2021density,achiam2017cpo,liu2022cvpo,kim2023sdac}; and policy-conditioned or compositional world models adapt predictions across policies or structured environments \citep{chen2024pcm,zhao2022howm,wang2025wm3c}. Ensembles and selective prediction estimate when a predictor may fail \citep{lakshminarayanan2017ensembles,geifman2019selectivenet}. Yet none of these pieces by itself identifies what must be covered when a separately queryable safety operator rewrites actions, or when a prediction error changes an irreversible return decision.

Our first insight is that \emph{pair novelty is not equivalent to executed-interface extrapolation}. An unseen pair \((\pi,\Pi)\) need not be statistically novel at the primitive interface if its induced \((x,u^{\rm exec})\) occupancy is identified. Conversely, observing each policy and each safety operator somewhere does not identify a target pair that reaches an uncovered executed interface. Our second insight is decision-theoretic: global Resource-to-Go error and mission-critical error are different. Under a threshold ReturnManager, bounded estimation error changes a decision only near the return boundary, and underestimation in this band is the direct route to late commitment.

We therefore study executed-interface Resource-to-Go, not a nominal-action energy label. We establish a paired identification/non-identification result, give an assumption-explicit finite-horizon transport bound with a proper stochastic-shortest-path truncation term, and prove return-boundary stability. The practical framework compares direct temporal-difference prediction, policy-conditioned and executed-action world models, and epistemic ensembles under a common irreversible ReturnManager. A more structured reliability mechanism is retained only if these simpler models fail a predeclared gate.

Our contributions are:
\begin{itemize}[leftmargin=*,topsep=2pt,itemsep=1pt]
  \item a formal executed-interface object that separates unseen pair identity from the support actually needed to identify charger Resource-to-Go;
  \item identification, impossibility, long-horizon error, and return-boundary results whose assumptions explicitly separate coverage from smoothness;
  \item a falsifiable protocol linking prediction error to return decisions through matched-support, matched-horizon, risk--coverage, and paired battery-cycle evaluations; and
  \item \resultpending{EMPIRICAL CLAIMS} an evidence slot for testing whether interface extrapolation and hitting horizon predict Resource-to-Go failure, and whether reliable estimates improve the stranding--throughput frontier.
\end{itemize}
